\scanpage{61}
\mspage{1.}

\uline{Th.} Soit $X$ préschéma alg., sur $k$ alg. clos. Alors
\[
\begin{aligned}
\pi'_1(X)={}&\text{limite projective des groupes quotients discrets de }\pi_1(X)\text{ d'\unsure{ordres|ordre}}\\
&\text{\unsure{premiers|premier} à }p
\end{aligned}
\]
est top. de type fini.

Grâce à ``van Kampen'', on peut supposer $X$ affine \inserted{\deleted{\unsure{intègre}} réduit}, puis $X$ affine et \uline{normal} [sinon, on aura
\diagram{\begin{tikzcd}[ampersand replacement=\&,column sep=small]
X'\times_X X'\arrow[r,shift left=.35ex]\arrow[r,shift right=.35ex] \& X'\arrow[r] \& X
\end{tikzcd}}
, et $\pi_1(X)$ est engendré par $\pi_1(X')$ et un nb fini de générateurs \unsure{correspondant|correspondants} aux composantes irréd. de $X'\times_X X'$, donc il en est de $=$ de $\pi'_1(X)$ \inserted{par $\pi'_1(X')$ \unsure{etc\ldots}}, au \deleted{\unsure{a}} sens des groupes \unsure{du|de} type \unsure{prodiscrets|profinis|profini}, \unsure{mot ops.|mut. ops.|mot $(p)$}

Si $\dim X=0$, trivial

Si $\dim X=1$, théorie de la spécialisation du groupe fondamental \inserted{``\unsure{tame|doux}''} nous ramène au cas de la caractéristique $0$, où on a \deleted{\unsure{un}} le \deleted{\unsure{groupe}} résultat par voie transcendante.

Si $\dim X\geq 2$, on considère une immersion
\diagram{\begin{tikzcd}[ampersand replacement=\&]
X^n\arrow[r] \& \unsure{\mathbb P|P}_k^r
\end{tikzcd}}
et un hyperplan \deleted{générique, $H$} \inserted{\unsure{variété|sous-espace} lin gén $L^{r-n+1}$}, défini sur une ext. alg close $K$ de $k$, on pose
\[
 C=\deleted{\unsure{H_{K}}}\,X^n\times_{\mathbb P_k^r}\deleted{\unsure{H}}L^{r-n+1}
\]
qui est une courbe normale de dim \unsure{un|$n$} sur $K$, et alors en vertu de Bertini ``générique'',
\diagram{\begin{tikzcd}[ampersand replacement=\&]
\pi_1(C)\arrow[r] \& \pi_1(X)
\end{tikzcd}}
est \uline{surjectif}. [\unsure{Réf|Ref}]. Donc $\pi'_1(C)\unsure{-|\longrightarrow}\pi'_1(X)$ l'est aussi, O.K.

\uline{Question} $X$ normale complète, $Y\subset X$ \deleted{\unsure{de codim} $\unsure{\geq 1|>1}$,} Alors $\pi_1(X-Y)$ mod le sous-gpe invariant fermé engendré par les groupes d'inertie wild en les composantes de \deleted{\unsure{degré}} codim 1 de $Y$, est de top. de type fini ??
% archive sheet number: 60
